
\section{Round-fourteen positive closure: physical impact suspension and a nondegenerate flag-current complex}
\label{sec:r14-a1}

The physical section and the symbolic natural extension are different objects.
The section is the unit square with its ordinary area form and with the seam
orbits removed; the two-sided shift is only a coding extension.  This
distinction removes the inverse-limit dimension error.  The response space
below is based on geometric mass on the fixed square, rather than an infimum
of cellwise norms with a depth-decaying factor.

\subsection{The coupled natural extension}

Let \(I_i(a)=(s_{i-1}(a),s_i(a))\), with width \(w_i(a)\), and put
\[
 B_a(q,p)=\left({q-s_{i-1}(a)\over w_i(a)},
                s_{i-1}(a)+w_i(a)p\right),
 \qquad q\in I_i(a).
\]
Let \(\Sigma^+=\{1,2,3,4\}^{\mathbb N_0}\) and
\(\Sigma^-=\{1,2,3,4\}^{\mathbb N}\).  For
\(x=(x_0,x_1,\ldots)\) and \(y=(y_1,y_2,\ldots)\), set
\[
 {\mathcal T}(x,y)=(\sigma x,x_0y),
 \qquad x_0y=(x_0,y_1,y_2,\ldots).
\]

\begin{theorem}[Physical coding and the coupled extension]
\label{thm:r14-a1-coding}
Let \(\mathcal S_a\) be the union of all forward and backward iterates of
the vertical and horizontal seams, and put
\(S_a^\circ=(0,1)^2\setminus\mathcal S_a\).  There is a bimeasurable coding
\[
 \mathfrak c_a:S_a^\circ\longrightarrow
 \Sigma^+\times\Sigma^-
\]
onto a full Bernoulli-measure set such that
\[
 \mathfrak c_a B_a={\mathcal T}\mathfrak c_a,
 \qquad
 (\mathfrak c_a)_\#(dq\,dp)=\nu_a^+\otimes\nu_a^- .
\]
In particular the newly created past symbol is the present future symbol;
there is no inverse of an independent one-sided past shift.
\end{theorem}

\begin{proof}
For a finite future word \(i_0\ldots i_n\), the corresponding horizontal
coordinate interval has length \(\prod_{k=0}^n w_{i_k}(a)\).  For a finite
past word \(j_1\ldots j_m\), the vertical coordinate interval has length
\(\prod_{\ell=1}^m w_{j_\ell}(a)\).  Their product rectangles are nested and
their diameters tend to zero.  Hence every two-sided word outside the endpoint
ambiguities determines one point, and every point outside \(\mathcal S_a\)
determines one word.  The rectangle areas factor into the Bernoulli weights.
Applying \(B_a\) deletes \(i_0\) from the future and prepends it to the past,
which is exactly \({\mathcal T}\).
\end{proof}

\subsection{An autonomous Hamiltonian impact suspension}

Write \(C_i=I_i(a)\times(0,1)\).  On \(C_i\) the branch \(B_{a,i}\) is an
exact symplectic diffeomorphism from the vertical rectangle \(C_i\) to the
horizontal rectangle \((0,1)\times I_i(a)\).  Its primitive difference is
\[
 B_{a,i}^*(p\,dq)-p\,dq=dG_{a,i},
 \qquad
 G_{a,i}(q,p)=
 {s_{i-1}(a)\over w_i(a)}(q-s_{i-1}(a))
 -s_{i-1}(a)p .
\]
We use open gates; their boundary is the impact singular set.

\begin{lemma}[Exact channel realization]
\label{lem:r14-a1-channel}
For each \(i\) there is an exact symplectic cobordism
\((\mathcal C_{a,i},d\lambda_{a,i})\), with disjoint incoming and outgoing
faces identified respectively with \(C_i\) and \(B_{a,i}(C_i)\), and an
autonomous Hamiltonian \(H_{a,i}\) whose regular incoming-to-outgoing
scattering map is \(B_{a,i}\).  The construction is smooth in \(a\) on every
compact subinterval of \((-1,1)\).
\end{lemma}

\begin{proof}
Use the isotopy
\[
 (q,p)\mapsto
 \left(s_{i-1}+e^{t\ell_i}(q-s_{i-1}),
       s_{i-1}+e^{-t\ell_i}(p-s_{i-1})\right),
 \qquad \ell_i=-\log w_i,
\]
and compose it with the two canonical translations needed to obtain the
displayed affine branch.  The three stages are Hamiltonian flows generated by
a linear momentum, by \(-\ell_i(q-s_{i-1})(p-s_{i-1})\), and by a linear
position Hamiltonian.  Choose a smooth clock partition
\(\chi_1+\chi_2+\chi_3=1\) on \(0<t<1\), so the stages have disjoint clock
support.  Adjoin a canonical clock pair \((t,E)\) and take
\(H=E+\sum_k\chi_k(t)H_k(q,p)\).  Then \(\dot t=1\), and the time-one
scattering map is exactly the composition, hence \(B_{a,i}\).  The open gate
\(C_i\) selects the channel geometrically; no external reset is applied.
The primitive \(G_{a,i}\) supplies the exact face matching of
\(\lambda=p\,dq+E\,dt-dG_{a,i}\) at the outgoing end.  Smooth dependence
follows from \(w_i(a)>0\).
\end{proof}

Let \(\mathcal M_a^\circ\) be the quotient of the disjoint channels by their
outgoing-to-incoming face identifications.  Only regular faces are glued;
seams are retained as ideal impact faces in the graph completion.

\begin{theorem}[Autonomous impact network with mechanical work]
\label{thm:r14-a1-impact}
The regular quotient \(\mathcal M_a^\circ\) is a Hausdorff exact symplectic
manifold carrying a complete autonomous Hamiltonian flow whose first return
to the physical square \(S_a^\circ\) is \(B_a\).  After adjoining a canonical
work pair \((z,\pi_z)\), the Hamiltonian
\[
 \widetilde H_a=H_a+\pi_z c_i\quad\hbox{on channel }i
\]
satisfies \(\dot z=c_i\), \(\dot\pi_z=0\).  Thus the increment of the genuine
canonical coordinate \(z\) in one return is \(c_i\), and after \(n\) returns
it is the additive current \(A_n\).
\end{theorem}

\begin{proof}
Every gluing map is the exact symplectic scattering map of
Lemma~\ref{lem:r14-a1-channel}; exact face matching gives a well-defined
primitive on the quotient.  Because the gates are disjoint open sets and the
gluing is between paired regular faces, local quotient charts are ordinary
mapping-torus charts.  Distinct gates remain separated, so the quotient is
Hausdorff.  The clock vector field crosses each face transversely, and its
return map is the channel scattering map.  The added term is Hamiltonian on
the product symplectic form
\(d\lambda_a+dz\wedge d\pi_z\), and Hamilton's equation gives the stated work
increment.
\end{proof}

\subsection{The graded geometric flag complex}

For integers \(m,n\ge0\), let \(\mathcal P_{m,n}(a)\) be the finite
polyhedral partition of the square generated by seam images through time
\(m\) and seam preimages through time \(n\).  A flag
\(F=(F_0\supset F_1\supset\cdots\supset F_k)\) is an oriented chain of faces
of codimension \(0,\ldots,k\).  Let \(\mathbf M(T)\) denote the mass of a
finite current \(T\).

For \(q\ge0\), define \(\mathscr F_a^q\) as the completion of compatible
families of currents \(T_F\), over all finite partitions and all
\(k\le q\), for the norm
\[
 \|T\|_{\mathscr F_a^q}
 =
 \sum_{k=0}^q \gamma^k
 \sup_{m,n}
 \sum_{\substack{F\in\mathcal P_{m,n}(a)\\\operatorname{codim}F=k}}
 \sum_{j=0}^{q-k}\mathbf M(D_a^jT_F),
 \qquad \gamma>0.
\]
Compatibility means restriction under refinement and cancellation of the two
opposite orientations on every internal face.  There is no factor decaying
with refinement depth.

\begin{proposition}[Nondegenerate complete flag complex]
\label{prop:r14-a1-flag}
The displayed expression is a norm.  Compatible refinement maps are
isometries on bulk mass and contractions on total flag mass.  The resulting
space is Hausdorff and complete, and the constant bulk density has norm one.
For each finite \(q\), material differentiation is bounded
\[
 \mathfrak D_a:\mathscr F_a^{q+1}\longrightarrow\mathscr F_a^q.
\]
The projective intersection
\(\mathscr F_a^\infty=\bigcap_q\mathscr F_a^q\) supports derivatives of every
finite order.
\end{proposition}

\begin{proof}
For a current restricted to a disjoint refinement, total variation is
additive; hence the sum of bulk masses is independent of the partition.  The
same statement holds on each oriented face after internal orientations are
identified.  Therefore a nonzero compatible current has positive mass on
some flag and cannot have zero norm.  In particular the restrictions of the
constant density have masses equal to the cell areas, whose sum is one.

Differentiating a moving \(k\)-face gives its material derivative on that face
plus the boundary contraction with its normal velocity, which is a
\((k+1)\)-flag current.  Thus one derivative raises codimension by at most one
and consumes one \(a\)-derivative.  The mass estimate for pushforward by the
affine branch and the finite seam velocity bound give the displayed bounded
map.  Completeness follows componentwise from completeness of finite-mass
currents and the closed compatibility relations.  Iterating at finite order
and taking the projective intersection proves the last assertion.
\end{proof}

\subsection{Typed response}

A transported observable is part of the data: it is a family
\(F_a\in\mathscr F_a^\infty\) obtained from one reference current by the
declared physical transport.  Let \(\ell_a(F)=\int F\,dq\,dp\).
The future and past Bernoulli conditional expectations have one-sided
reduced resolvents \(R_a^\pm\) on their centered Hölder spaces.

\begin{theorem}[Bulk--score--flag response identity]
\label{thm:r14-a1-response}
For every \(q\) and every \(C^{q+1}\) transported family \(F_a\),
\[
 {d^q\over da^q}\ell_a(F_a)
\]
is a finite sum of words formed from: the one-sided reduced resolvents
\(R_a^\pm\), derivatives of the Bernoulli score, the coupled present-symbol
insertion \(x_0y\), and the flag derivatives
\(\mathfrak D_a^jF_a\), with every word bounded
\(\mathscr F_a^{q+1}\to\mathbb R\).  The formula is independent of finite
cylinder presentation.  If \(F_a\equiv F\) is a genuinely fixed physical
observable, then all score and seam terms cancel and the response is zero.
\end{theorem}

\begin{proof}
First take a finite bi-cylinder observable.  Differentiate its exact rectangle
integral.  Interior differentiation gives the transported bulk term,
differentiation of Bernoulli weights gives the centered score, and the
Reynolds boundary formula gives the oriented seam current.  Conditional
martingale decomposition of the score into future and past coordinates
produces the two one-sided Poisson resolvents; the present symbol is inserted
through the coupled map of Theorem~\ref{thm:r14-a1-coding}, not through two
independent shifts.  Repeated differentiation yields flags of codimension at
most \(q\), exactly as in Proposition~\ref{prop:r14-a1-flag}.

The mass norm makes common refinement an isometry, so the formula agrees on
all presentations and extends by density.  For a fixed physical \(F\),
\(\ell_a(F)=\int_{(0,1)^2}F\,dq\,dp\) is independent of \(a\); the differentiated
change-of-variables identity gives the asserted cancellation.
\end{proof}

\begin{theorem}[Round-fourteen A1 closure]
\label{thm:r14-a1-main}
The generalized-baker benchmark has a correctly coupled symbolic natural
extension, an autonomous exact Hamiltonian impact realization with a
canonical work coordinate, and a nondegenerate graded flag-current response
calculus of every finite order.
\end{theorem}

\begin{proof}
Combine Theorems~\ref{thm:r14-a1-coding},
\ref{thm:r14-a1-impact}, and \ref{thm:r14-a1-response}, with
Proposition~\ref{prop:r14-a1-flag}.
\end{proof}
